% Proposed insertion into Applications; no unrelated manuscript text modified.
\mypara{Visual Concepts Associated with Human Preference}
We apply dataset diffing within recorded human judgments, contrasting preferred images with less-preferred alternatives under a matched prompt. Frozen FG-CLIP2 scores use the maximum normalized text--patch similarity over valid patches and the full supplied bank of 54,030 concepts. This pilot uses a 1,024-patch budget. Similarity scores are not calibrated concept-presence probabilities.

In ImageRewardDB, lower human rank defines preference. We average all strict ranked comparisons within each annotation event, events within each normalized prompt, and then prompts. Official training data selects the ten most positive and ten most negative concepts; identities and directions are frozen before official-test evaluation. Prompt/image dependencies crossing official split boundaries are quarantined, and incomplete events are excluded. A 500-prompt sample leaves 337 discovery, 75 validation and 75 test prompts after image decoding failures. The test set contains 1,043 strict comparisons from 75 independent prompt/image components.

An equal signed mean of the twenty frozen concept scores achieves prompt-balanced paired concordance 0.5513 (95\% cluster-bootstrap interval [0.4993, 0.5993], 2,000 resamples), counting ties as half. The preferred-minus-rejected mean score difference is +0.002243 [0.000622, 0.003845]. A positive mean margin and uncertain ordering performance are compatible because these measure different aspects of the score distribution. Nineteen of twenty concepts retain their discovery effect direction, without a simultaneous significance claim. Preferred-side candidates include \textit{lip-syncing}, \textit{stare} and \textit{smirk}; rejected-side candidates include \textit{tarmac}, \textit{pavement} and \textit{asphalt}. These lexical associations remain human-unverified. Patch maps and observed paired examples expose possible mismatches between a concept's wording and its image activations.

Best-versus-rest and best-versus-worst sensitivities reuse image scores but independently freeze their candidates on discovery data. Their concordances are respectively 0.4840 [0.4173, 0.5535] and 0.5967 [0.4900, 0.6967]. All three intervals include 0.5. The results therefore do not establish reliable preference ordering across constructions. Generator and annotator identities are unavailable in the ImageReward metadata; prompt matching does not remove all confounding or establish causality. The pilot overlaps an earlier smaller engineering run and is not an independent replication.

The linked four-candidate Pick-a-Pic tables have 13,085 strictly eligible prompt-matched events, but the selected original image URLs failed, and additional UID-verifiable archives yielded no complete eligible event. Empirical Pick-a-Pic effects and cross-dataset replication are \emph{unavailable}; we do not substitute a different release's annotations. An independent text-bank diagnostic supports templated box-mode compatibility (mean cosine 0.9656 on 32 evenly spaced entries), while the exact bank-generation recipe remains unconfirmed. Human concept verification, restored original images and larger independent evaluation are still needed.
